\begin{figure}[H]
\centering
\begin{tikzpicture}[
  font=\footnotesize,
  actor/.style={draw=blue!55!black, rounded corners=2pt, fill=blue!5,
    align=center, text width=2.65cm, minimum height=.8cm, inner sep=4pt},
  message/.style={-{Latex[length=2mm]}, thick, draw=blue!65!black},
  ack/.style={-{Latex[length=2mm]}, thick, draw=green!45!black}]
\node[actor] at (0,0) {Simulator\\publisher MQTT};
\node[actor] at (3.75,0) {Mosquitto\\broker MQTT};
\node[actor] at (7.5,0) {Bridge\\subscriber + producer};
\node[actor] at (11.25,0) {Kafka\\broker};
\foreach \x in {0,3.75,7.5,11.25} {\draw[dashed, black!40] (\x,-.5) -- (\x,-6.1);}
\draw[message] (0,-1.1) -- node[above]{1. PUBLISH QoS 1} (3.75,-1.1);
\draw[ack] (3.75,-1.95) -- node[above]{2. PUBACK cho publisher} (0,-1.95);
\draw[message] (3.75,-2.8) -- node[above]{3. PUBLISH QoS 1} (7.5,-2.8);
\draw[message] (7.5,-3.65) -- node[above]{4. produce, \texttt{acks=all}} (11.25,-3.65);
\draw[ack] (11.25,-4.5) -- node[above]{5. delivery callback: OK} (7.5,-4.5);
\draw[ack] (7.5,-5.35) -- node[above]{6. PUBACK thủ công} (3.75,-5.35);
\node[draw=orange!60!black, fill=orange!7, rounded corners=2pt,
  text width=13.3cm, align=left, inner sep=5pt] at (5.625,-6.95)
  {\textbf{Khoảng có thể phát lại:} Bridge ngừng sau khi Kafka nhận bản tin nhưng trước bước 6.
   Mosquitto giữ bản chưa ACK trong phiên bền và có thể phát lại khi Bridge kết nối lại;
   định danh \texttt{event\_id} giúp hạ lưu nhận diện bản trùng.};
\end{tikzpicture}
\caption{Hai vòng xác nhận MQTT QoS 1 độc lập và điểm ACK thủ công của Bridge. Đây là một trình tự minh hoạ; PUBACK cho publisher không xác nhận rằng Kafka hoặc PostgreSQL đã ghi thành công.}
\label{fig:mqtt-ack-boundary}
\end{figure}
